\documentclass[11pt,twoside]{amsart}
\usepackage{amssymb, amsmath, enumerate, palatino, hyperref}
\usepackage[normalem]{ulem}
\usepackage{fullpage}
\usepackage[T1]{fontenc}
\renewcommand{\labelitemi}{\guillemotright}
\usepackage{mathrsfs}

\usepackage{tikz}
\tikzstyle{ball} = [circle,shading=ball, ball color=black,
    minimum size=1mm,inner sep=1.3pt]

\theoremstyle{plain}
\newtheorem{prop}{Proposition}%[section]
\newtheorem{lemma}[prop]{Lemma}
\newtheorem{thm}[prop]{Theorem}
\newtheorem{obs}[prop]{Observation}
\newtheorem{app}[prop]{Application}
\newtheorem*{MainThm}{Main Theorem}
\newtheorem{cor}[prop]{Corollary}
\newtheorem{conj}[prop]{Conjecture}
\theoremstyle{remark}
\newtheorem{rmk}[prop]{Remark}
\newtheorem{prob}{Problem}
\newtheorem{bonus}[prop]{Bonus Problem}
\newtheorem{exc}{Exercise}
\theoremstyle{definition}
\newtheorem{ex}[prop]{Example}
\theoremstyle{definition}
\newtheorem{defn}[prop]{Definition}

\newcommand{\RR}{\mathbb{R}}
\newcommand{\ZZ}{\mathbb{Z}}
\newcommand{\CC}{\mathbb{C}}
\newcommand{\NN}{\mathbb{N}}
\newcommand{\QQ}{\mathbb{Q}}
\newcommand{\PP}{\mathbb{P}}

\newcommand{\cC}{\mathscr{C}}
\newcommand{\Ob}{\operatorname{Ob}}
\newcommand{\Mor}{\operatorname{Mor}}

\newcommand{\Set}{\operatorname{Set}}
\newcommand{\FinSet}{\operatorname{FinSet}}
\newcommand{\Vect}{\operatorname{Vect}}
\newcommand{\FinVect}{\operatorname{FinVect}}
\newcommand{\Mat}{\operatorname{Mat}}
\newcommand{\Gp}{\operatorname{Gp}}
\newcommand{\FinGp}{\operatorname{FinGp}}
\newcommand{\AbGp}{\operatorname{AbGp}}
\newcommand{\Ring}{\operatorname{Ring}}
\newcommand{\CommRing}{\operatorname{CommRing}}
\newcommand{\Field}{\operatorname{Field}}
\newcommand{\Top}{\operatorname{Top}}
\newcommand{\Aut}{\operatorname{Aut}}
\newcommand{\id}{\operatorname{id}}
\DeclareRobustCommand{\stir}{\genfrac\{\}{0pt}{}}


\newcommand{\defeq}{:=}

\title{Math 113: Discrete Structures\\ Homework due Wednesday Week 10}
%\author{} %Remove the leading % and put your name here if using this .tex file as a template for your homework.

\begin{document}
\maketitle

\noindent
\textbf{Due:} Wednesday, April 8 at 10pm.
\bigskip

\begin{prob}
  You have three coins.  Two of the coins are fair: when flipped they are
  equally likely to land heads or tails. One coin, however, is weighted somehow
  so that its probability of landing heads is~$2/3$.
  \begin{enumerate}[(a)]
    \item Choose one of the three coins uniformly at random and flip it. What is
      the probability the result is heads?  For your solution, number the
      coins~$1,2,3$ with coin~$3$ being the weighted one, and let~$A_i$ denote
      the event that coin~$i$ was chosen. Apply the generalized law of total
      probability.
    \item Choose one of the three coins at random and flip it.  It lands heads.
      What is the probability that you chose the weighted coin? (Hint: Bayes'
      law.)
     \item Choose one of the three coins at random and flip it twice.  Let $E$ be the event that you get heads on the first flip and $F$ be the even that you get heads in the second flip. Are $E $ and $F$ independent? Prove your answer.
  \end{enumerate}
\end{prob}

\begin{prob}
  (We reconsider one of the challenge problems from an in-class worksheet.)

 Your friend invites you to play a game: they write ten distinct real numbers on
ten blank cards.  The cards are shuffled randomly and placed face down on the
table.  You start at the top of the deck and start revealing cards.  At any
point you may choose to stop turning over cards and select the most recently
revealed card.  You win if your selection is the largest of all ten numbers
(both those previously revealed and those still unrevealed). 

It is beyond the scope of this class to prove that the best strategy is given by the \emph{stopping rule}. For the stopping rule, we fix a number $r$ such that $1\leq r \leq 10$. We look at the first $r-1$ cards and note the maximal value among them, $M$.  For the subsequent $10-(r-1)$ cards, select the first one larger than $M$. We will compute the probability of winning when using this strategy. For the remainder of the problem, we will assume $r>1$ (the case $r=1$ can be analyzed directly). 
\begin{enumerate}[(a)]
 \item For $1\leq a \leq 10$, assume that the number in position $a$ is the largest in the deck. What is the probability that we select $a$ with this strategy? (Note that your answer will depend on whether or not $a<r$.)
 \item Use the law of total probability and your answer to (a) to compute the probability that you win using this strategy.
 \item Use a calculator or computer to compute the probabilities of winning with this strategy for all $2\leq r \leq 10$. Which $r$ gives you the best chances of winning?
\end{enumerate}
\end{prob}
\end{document}
